\chapter{Arithmetic Hilbert--P\'olya and the VB\(^{*}\) Converse: Structural Analyses}
\label{v4-arith:HPAr-vbconverse}

\emph{Assume P1, P2, and P3 in the comparison form used in
Chapter~\ref{v4-arith:rh-reformulation}.  The two additional arithmetic
inputs are H-P\(^{\mathrm{Ar}}\) and VB\(^{*}\).  H-P\(^{\mathrm{Ar}}\)
is the determinant-trace hypothesis: there is a self-adjoint operator
\(A\) for which
\(\Theta_{\xi}=1/2+iA\) has determinant divisor
\(\mathrm{div}\,\xi_{\mathbb R}\) and Weil--Connes trace equal to the
zero distribution.  VB\(^{*}\) adds positivity and Hamburger
determinacy for the Fock-space moment measure.  Connes' adelic trace
formula and Deninger's regularised determinant give equivalent forms of
the first input after the named intertwiners; de Branges positivity and
Nyman--Beurling--B\'aez-Duarte decay give equivalent forms of the
second.  Zero-placement supplies support; it does not supply the
unbounded moment theorem.}

\section{The Analytic Splitting}
\label{v4-arith:hpvbconv:context}

Chapter~\ref{v4-arith:rh-reformulation} gives a sufficient implication
to RH from four inputs: arithmetic Positselski, Koszul--Petersson
compatibility, H-P\(^{\mathrm{Ar}}\), and VB\(^{*}\).  The P1 input is
the completed arithmetic Positselski comparison; its Tate-Poisson
shadow is only the forward functional equation.  The P3 input is the
full Koszul--Petersson comparison; its proved sph-fin cuspidal core is
only the Tamagawa normalisation.  The last two inputs are the
Hilbert--P\'olya determinant-trace realisation and the Fock-space
positivity-determinacy theorem.

\subsection{The Asymmetry}
\label{v4-arith:hpvbconv:gap-structure}

Zero-placement and VB\(^{*}\) have different strength. Three
operator-theoretic facts account for the difference:

\begin{enumerate}
\item Moment data for an unbounded operator do not determine the
spectral measure without a determinacy theorem. Akhiezer--Glazman
Ch.~VI.3 supplies the classical obstruction.
\item The assertion that the normal zero-divisor operator has determinant
divisor \(\mathrm{div}\,\xi_{\mathbb R}\) and Weil--Connes zero trace
is precisely H-P\(^{\mathrm{Ar}}\)
(\Cref{v4-arith:rh:hyp:HPAr}).
\item A symmetric \(L_{0}^{\mathrm{Ar}}\) on the primary Fock is not
automatically a self-adjoint Hilbert--P\'olya operator; a domain and a
self-adjoint extension must be specified.
\end{enumerate}

Thus \Cref{v4-arith:rh:thm:VB-implies-F-RH} gives the implication
VB\(^{*}\Rightarrow\) zero-placement under P2, P3, and
H-P\(^{\mathrm{Ar}}\). The converse implication needs spectral
regularity: positivity of the regularised moment functional,
Hamburger determinacy, and support compatibility for the Fock
operator.

\subsection{Disjoint Sources}
\label{v4-arith:hpvbconv:sources}

For H-P\(^{\mathrm{Ar}}\):
\begin{description}
\item[Connes, arXiv:math/9811068 (1999).] Trace formula on the adele
class space \(X_{\mathbb{Q}}=\mathbb{A}_{\mathbb{Q}}/\mathbb{Q}^{\times}\);
absorption trace of the scaling action and the zero distribution of
\(\zeta\). Selecta~Math.~5 version with full technical details.
\item[Connes--Consani--Marcolli 2007.] Adele
class space with quantum-statistical-mechanical structure;
KMS\({}_{\beta=1}\) symmetry breaking.
\item[Deninger, Proc.~Symp.~Pure Math.~55.1 (1994).] Motivic
\(L\)-functions; formalism identifying \(\xi(s)\) with a regularised
determinant of a Frobenius operator on an arithmetic cohomology.
\item[Deninger, Math.~Nachr.~205 (2001).] Motivic cohomology
dynamical systems formalism.
\item[Berry--Keating, SIAM Review 41 (1999).] Proposal that
\(H=\tfrac{1}{2}(xp+px)\) on \(L^{2}(\mathbb{R}_{>0},dx/x)\) is the
Hilbert--P\'olya candidate; restriction to a Planck-cell cutoff gives the
correct leading asymptotic zero density.
\item[Akhiezer--Glazman \emph{Theory of Linear Operators in Hilbert
Space} Vol.~II (Dover 1993 reprint).] Ch.~VI on unbounded self-adjoint
operators, Friedrichs extensions, and moment problems; Ch.~VI.3 is the
Hamburger/Carleman determinacy criterion.
\item[Reed--Simon \emph{Methods of Modern Mathematical Physics} Vol.~I
(Academic Press 1980), Vol.~II (1975).] Ch.~VIII (Vol.~I) on
unbounded operators; Vol.~II on Fourier analysis and
self-adjointness.
\end{description}

For the VB\(^{*}\) converse:
\begin{description}
\item[de Branges \emph{Hilbert Spaces of Entire Functions} (Springer
1968).] Foundational theorem on reproducing-kernel Hilbert spaces of
entire functions; Hermite--Biehler structure theorem.
\item[Lagarias, Acta Arith.~89 (1999), 217--234.] Reformulation of
RH as Hermite--Biehler positivity for
\(E_{\xi}(z)=\xi_{\mathbb R}(\tfrac{1}{2}-iz)+
\xi_{\mathbb R}'(\tfrac{1}{2}-iz)\).
\item[B\'aez-Duarte, Adv.~Math.~190 (2005), 315--320; arXiv:math/0202141.]
Quantitative strengthening of the Nyman--Beurling criterion.
\item[Nyman \emph{On Some Groups and Semigroups of Translations}
(Uppsala 1950).] Original criterion: RH equivalent to density of
\(\{\{\theta/x\}: \theta\in(0,1]\}\) translates in
\(L^{2}((0,1))\).
\item[Beurling, Proc.~Natl.~Acad.~Sci.~41 (1955).] Beurling sequence
space formulation of Nyman.
\item[Burnol 2000.] Hermite--Biehler
refinement of Nyman--Beurling.
\item[Simon, Adv.~Math.~137 (1998), 82--203.] Classical moment problem
as self-adjoint finite difference operator; unbounded case treated via
Jacobi matrix formalism.
\item[Stein--Weiss \emph{Introduction to Fourier Analysis on
Euclidean Spaces} (Princeton 1971).] Ch.~V on Hardy spaces and
boundary-value theorems governing Hermite--Biehler.
\end{description}

These analytic sources supply the external operator models for
H-P\(^{\mathrm{Ar}}\) and the VB\(^{*}\) converse.  P1, P2, and P3
enter only through their comparison hypotheses.

\section{Structural Reductions}
\label{v4-arith:hpvbconv:table}

The assertions below separate domain, trace, determinant, and moment
questions. Each assertion is an implication between named mathematical
structures.

\subsection{H-P\(^{\mathrm{Ar}}\): Operator-Theoretic Framework}
\label{v4-arith:hpvbconv:cycle-1}

On the finite-span core of the primary arithmetic Heisenberg Fock, the
Sugawara operator is diagonal with real non-negative eigenvalues
(\Cref{v4-arith:vb:lem:diagonal-action}). Reed--Simon Vol.~I
Thm.~VIII.11 and the Friedrichs construction give a canonical
self-adjoint non-negative extension from this semibounded quadratic
form. This removes the domain ambiguity in
\Cref{v4-arith:rh:hyp:HPAr}. It does not identify the resulting
Bost--Connes partition trace with the Weil--Connes zero distribution;
that identification is the spectral trace input of
\Cref{v4-arith:hpvbconv:prop:BC-trace-compatibility}.

\subsection{H-P\(^{\mathrm{Ar}}\): Connes Adele Class Space}
\label{v4-arith:hpvbconv:cycle-2}

Connes' adelic trace formula is a theorem on
\(X_{\mathbb{Q}}=\mathbb{A}_{\mathbb{Q}}/\mathbb{Q}^{\times}\). The
comparison with H-P\(^{\mathrm{Ar}}\) requires an intertwiner from the
primary Fock space to the adele-class \(L^{2}\)-space which carries the
zero-mode Bost--Connes monoid action to the idele-class scaling
action. Under such an intertwiner, Connes' absorption trace and the
regularised Fock trace are the same distribution. The displayed
absorption-trace formula,
\[
\sum_{\rho}h(\mathrm{Im}\,\rho)
=h(0)+h(1)-\hat{h}(0)\log\pi-\sum_{v}\int_{\mathbb{Q}_{v}^{\times}}'\frac{h(\log|u|_{v})}{|1-u|_{v}}d^{\times}u,
\]
matches (as a distribution on the dual real line) the Li--Stieltjes
moment generating function
\(\sum_{\rho}e^{-t\rho}\) through Mellin inversion
\(h(u)=e^{-tu/2}\), provided the embedding preserves the
distributional trace.

\subsection{H-P\(^{\mathrm{Ar}}\): Deninger Motivic Cohomology}
\label{v4-arith:hpvbconv:cycle-3}

Deninger's formalism postulates a cohomology
\(H^{\bullet}_{M}(\overline{\mathrm{Spec}\,\mathbb{Z}},\mathbb{R}(\cdot))\)
with Frobenius \(\Phi\) and Lefschetz determinant formula. If that
cohomology exists and is identified with the primary Fock by an
intertwiner, then the determinant identity for \(\xi_{\mathbb R}\) is
the H-P\(^{\mathrm{Ar}}\) determinant statement. The formal identity
\(\xi(s)=\det_{\infty}(s\cdot\mathrm{id}-L_{0}^{\mathrm{Ar}})\)
holds under H-P\(^{\mathrm{Ar}}\): the regularised determinant is defined by
\(\log\det_{\infty}(A)=-\partial_{s}\mathrm{tr}(A^{-s})|_{s=0}\), and
the determinant divisor is \(\mathrm{div}\,\xi_{\mathbb R}\) after the
trivial-zero and pole summands have been removed by the primary-sector
restriction.

\subsection{H-P\(^{\mathrm{Ar}}\): Necessary and Sufficient Conditions}
\label{v4-arith:hpvbconv:cycle-4}

The three formulations have one common mathematical core: a
positive-definite Hilbert space \(\mathcal{H}\), a densely-defined
self-adjoint operator \(A\), and an intertwiner
\(\pi:\mathcal{V}^{\mathrm{prim}}_{\mathrm{Heis},\mathrm{prim}}\to\mathcal{H}\)
such that \(A\pi=\pi L_{0}^{\mathrm{Ar}}\), the normal operator
\(\tfrac12+iA\) has regularised determinant divisor equal to the
non-trivial zero divisor of \(\xi_{\mathbb R}\), and its
Weil--Connes trace is the corresponding zero distribution. This is
classical Hilbert--P\'olya inside the arithmetic-Fock framework.

\subsection{VB\(^{*}\) Converse: Spectral Regularity Beyond RH}
\label{v4-arith:hpvbconv:cycle-5}

Zero-placement gives the zeros on the critical line. VB\(^{*}\) asks
for positivity of the operators \(T_{n}\) on the Fock closure. Passing
from the former to the latter is a Hamburger moment problem for an
unbounded operator; in this setting moment positivity does not select
a unique spectral measure. Carleman's criterion says that the
Hamburger
moment problem with moments \(\{\mu_{n}\}_{n\geq 0}\) is determinate if
\(\sum_{n\geq 1}\mu_{2n}^{-1/(2n)}=+\infty\). For the symmetrised zero
measure \(\nu\) of \Cref{v4-arith:vb:def:zero-measure},
\(\mu_{n}=\sum_{\rho}(\tfrac{1}{4}+\gamma_{\rho}^{2})^{n}\) in the
RH-conditional sense. Existing Riemann--von Mangoldt asymptotics do
not by themselves provide a rigorous positive regularised moment
measure with sharp two-sided \(\mu_{2n}\) bounds. The Carleman construction is
therefore open, not refuted.

\subsection{VB\(^{*}\) Converse: de Branges and B\'aez-Duarte}
\label{v4-arith:hpvbconv:cycle-6}

The de Branges and Nyman--Beurling formalisms give another expression
of the same missing regularity. RH is equivalent to the
Hermite--Biehler property of \(E_{\xi}\) in de Branges' sense, but
VB\(^{*}\) requires the quantitative B\'aez-Duarte completeness rate
compatible with the primary-Fock Petersson norm. B\'aez-Duarte
Thm.~1 of arXiv:math/0202141 gives
RH \(\Leftrightarrow\) the functions \(\{\rho_{n}(x)=\{n/x\}-n\{1/x\}\}_{n\geq 2}\)
span a dense subspace of \(L^{2}((0,1),dx)\) with a sharp decay rate
\(\epsilon_{N}=\mathrm{dist}(\mathbf{1},\mathrm{span}\{\rho_{n}\}_{n\leq
N})\) satisfying \(\epsilon_{N}\ll(\log N)^{-1/2}\) under RH and
\(\ll(\log N)^{-c}\) for some \(c>1/2\) only under additional spectral
regularity.

\section{H-P\(^{\mathrm{Ar}}\): Structural Analysis}
\label{v4-arith:hpvbconv:HPAr-analysis}

\subsection{Operator-Theoretic Framework}
\label{v4-arith:hpvbconv:operator-framework}

\begin{definition}[candidate \(L_{0}^{\mathrm{Ar}}\) with explicit domain]
\label{v4-arith:hpvbconv:def:L0-domain}
\ClaimStatusDefinitional. On the primary arithmetic Heisenberg Fock
\(\mathcal{V}^{\mathrm{prim}}_{\mathrm{Heis},\mathrm{prim}}\), equipped
with the Petersson inner product \(\langle\cdot,\cdot\rangle^{\mathrm{Pet}}\)
of \Cref{v4-arith:thm:P3-resolution}, define:
\begin{itemize}
\item \emph{Finite-span core}
\(\mathcal{D}_{0}\;:=\;\mathrm{span}_{\mathbb{C}}\{|\mathbf{n}\rangle\otimes\delta_{m}\;:\;\mathbf{n}\in\mathbb{Z}_{\geq
0}^{\oplus\mathbb{N}},\;m\in\mathbb{N}^{\times}\}\),
the linear span of number-state/zero-mode-lattice basis vectors
(\Cref{v4-arith:vb:lem:diagonal-action},
\Cref{v4-arith:def:zero-mode-lattice}).
\item \emph{Candidate operator}
\(L_{0}^{\mathrm{Ar},\,\mathrm{symm}}:\mathcal{D}_{0}\to\mathcal{V}^{\mathrm{prim}}_{\mathrm{Heis},\mathrm{prim}}\)
defined by
\begin{equation}
\label{v4-arith:hpvbconv:eq:L0-def}
L_{0}^{\mathrm{Ar},\,\mathrm{symm}}\bigl(|\mathbf{n}\rangle\otimes\delta_{m}\bigr)\;=\;\Bigl(\sum_{i\geq
1}i\,n_{i}\;+\;\lambda(m)\Bigr)|\mathbf{n}\rangle\otimes\delta_{m}
\end{equation}
where \(\lambda(m)\in\mathbb{R}\) is the Bost--Connes zero-mode weight
corresponding to \(m\in\mathbb{N}^{\times}\) (unspecified at this stage;
the choice is constrained by the Bost--Connes KMS\({}_{1}\) state
structure of \Cref{v4-arith:thm:BC-acts-on-L}).
\item \emph{Closure}: let \(L_{0}^{\mathrm{Ar}}\) denote the closure of
\(L_{0}^{\mathrm{Ar},\,\mathrm{symm}}\) with respect to the graph norm
on \(\mathcal{D}_{0}\).
\end{itemize}
\end{definition}

\begin{theorem}[Friedrichs self-adjoint extension]
\label{v4-arith:hpvbconv:thm:L0-Friedrichs}
\ClaimStatusProvedHereConditional. Under the hypothesis that \(\lambda(m)\geq 0\)
for all \(m\in\mathbb{N}^{\times}\) (non-negativity of the Bost--Connes
zero-mode weights), the symmetric operator
\(L_{0}^{\mathrm{Ar},\,\mathrm{symm}}\) of
\Cref{v4-arith:hpvbconv:def:L0-domain} admits a unique Friedrichs
self-adjoint extension
\(L_{0}^{\mathrm{Ar}}:\mathcal{D}\to\mathcal{V}^{\mathrm{prim}}_{\mathrm{Heis},\mathrm{prim}}\)
with \(\mathcal{D}\supset\mathcal{D}_{0}\) dense. On the Friedrichs
extension, \(L_{0}^{\mathrm{Ar}}\) is essentially self-adjoint on the
finite-span core \(\mathcal{D}_{0}\) and has non-negative spectrum.
\end{theorem}

\begin{proof}
Apply Reed--Simon Vol.~II Thm.~X.23 (Friedrichs extension of a
non-negative symmetric operator). The quadratic form
\(q(\psi,\psi')=\langle L_{0}^{\mathrm{Ar},\,\mathrm{symm}}\psi,\psi'\rangle^{\mathrm{Pet}}\)
is closable because \(L_{0}^{\mathrm{Ar},\,\mathrm{symm}}\) acts
diagonally on the orthonormal basis
\(\{|\mathbf{n}\rangle\otimes\delta_{m}\}\) with non-negative real
eigenvalues. Reed--Simon Vol.~I Thm.~VIII.11 then gives essential
self-adjointness on the finite-span core: a diagonal operator with
real eigenvalues on a Hilbert space with eigenvector basis is
essentially self-adjoint on the finite-span of eigenvectors. The
Friedrichs extension is the unique self-adjoint extension dominated
by the quadratic form (Akhiezer--Glazman Vol.~II Thm.~123). Non-negativity
of the spectrum follows from non-negativity of all eigenvalues.
\end{proof}

\begin{proposition}[Bost--Connes trace compatibility is the HP input]
\label{v4-arith:hpvbconv:prop:BC-trace-compatibility}
\ClaimStatusProvedHere. The Friedrichs extension of
\Cref{v4-arith:hpvbconv:thm:L0-Friedrichs} constructs a
self-adjoint non-negative Fock/BC Hamiltonian from the real weights
\(\lambda(m)\). It does not by itself construct the Hilbert--P\'olya
determinant-trace operator \(H_{\mathrm{HP}}\). H-P\(^{\mathrm{Ar}}\) holds only
after one supplies an additional distributional trace identification
with the zero side. The formal heat-kernel expression is
\begin{equation}
\label{v4-arith:hpvbconv:eq:BC-trace-condition}
\mathrm{tr}_{\mathrm{BC}}\bigl(e^{-tL_{0}^{\mathrm{Ar}}}\bigr)\;=\;\sum_{\rho}e^{-t\rho}
\qquad(t>0)
\end{equation}
but \eqref{v4-arith:hpvbconv:eq:BC-trace-condition} is only symbolic
notation. It means equality of Weil--Connes
distributions on admissible Paley--Wiener tests, equivalently the
trace identity of \Cref{v4-arith:rh:hyp:HPAr} for
\(\Theta_{\xi}=1/2+iH_{\mathrm{HP}}\). This trace compatibility is
the Hilbert--P\'olya determinant-trace input, not a consequence of the
Friedrichs theorem.
\end{proposition}

\begin{proof}
Direct computation of \(\mathrm{tr}(e^{-tL_{0}^{\mathrm{Ar}}})\) on the
basis of \(\mathcal{V}^{\mathrm{prim}}_{\mathrm{Heis},\mathrm{prim}}=L\otimes\mathcal{H}^{(\infty)}\):
\begin{equation*}
\mathrm{tr}\bigl(e^{-tL_{0}^{\mathrm{Ar}}}\bigr)\;=\;\Bigl(\sum_{m\geq
1}e^{-t\lambda(m)}\Bigr)\cdot\Bigl(\sum_{\mathbf{n}}e^{-t\sum_{i}i\,n_{i}}\Bigr)\;=\;\Bigl(\sum_{m\geq
1}e^{-t\lambda(m)}\Bigr)\cdot\prod_{i\geq 1}\frac{1}{1-e^{-ti}}.
\end{equation*}
For the Bost--Connes/KMS weight \(\lambda(m)=\log m\), the
zero-mode factor is \(\zeta(t)\) for \(\mathrm{Re}(t)>1\), multiplied
by the Dedekind-eta inverse from the Fock descendants. This is the
Euler/partition side. It is not the zero distribution. Passing from
this partition trace to the zero side requires the full
Weil--Connes explicit-formula distribution and a positive trace
realisation; that is exactly H-P\(^{\mathrm{Ar}}\) in
\Cref{v4-arith:rh:hyp:HPAr}. Thus the proposition records an
equivalence of formulations after the trace realisation is assumed,
not a derivation of the trace realisation.
\end{proof}

\begin{remark}[what the Friedrichs construction buys]
\label{v4-arith:hpvbconv:rem:friedrichs-gain}
\Cref{v4-arith:hpvbconv:thm:L0-Friedrichs} removes only one
domain-specification ambiguity: a real non-negative diagonal
Fock/BC Hamiltonian has a canonical self-adjoint extension. What
remains is not an elementary choice of \(\lambda(m)\). The KMS
normalisation selects \(\lambda(m)=\log m\), which gives the
Bost--Connes partition trace \(\zeta(t)\); the zero distribution of
\(\xi_{\mathbb{R}}\) (the completed xi function
\(\xi(s)=\tfrac{1}{2}s(s-1)\pi^{-s/2}\Gamma(s/2)\zeta(s)\),
here denoting the restriction to the primary archimedean sector) is a different object. The Hilbert--P\'olya
content is the additional construction of \(H_{\mathrm{HP}}\), or an
intertwiner carrying this Fock/BC Hamiltonian to the zero-divisor operator
\(\Theta_{\xi}\), with the positive Weil--Connes trace preserved.
\end{remark}

\subsection{Connection to Connes Adele Class Space}
\label{v4-arith:hpvbconv:connes-connection}

\begin{definition}[Connes adele class space]
\label{v4-arith:hpvbconv:def:adele-class-space}
\ClaimStatusDefinitional (Connes arXiv:math/9811068 \S II). The
\emph{adele class space} is
\(X_{\mathbb{Q}}=\mathbb{A}_{\mathbb{Q}}/\mathbb{Q}^{\times}\), the
quotient of the adeles by the multiplicative action of non-zero
rationals. Its measure-zero part \(X_{\mathbb{Q}}^{0}\) supports an
\(L^{2}\) space \(L^{2}_{0}(X_{\mathbb{Q}},d\mu)\) on which the idele
class group \(C_{\mathbb{Q}}=\mathbb{A}_{\mathbb{Q}}^{\times}/\mathbb{Q}^{\times}\)
acts unitarily by multiplication. The absorption trace of this action
is the content of Connes' realisation.
\end{definition}

\begin{theorem}[Connes arXiv:math/9811068 Thm.~1]
\label{v4-arith:hpvbconv:thm:connes-trace-formula}
\ClaimStatusProvedElsewhere. For appropriate test functions \(h\) on
\(C_{\mathbb{Q}}\),
\begin{equation}
\label{v4-arith:hpvbconv:eq:connes-trace-formula}
\mathrm{Tr}(\vartheta(h))\;=\;2h(1)+2h(-1)-\hat{h}(0)\log\pi-\sum_{v}\int_{\mathbb{Q}_{v}^{\times}}'\frac{h(u)}{|1-u|}d^{\times}u-\sum_{\rho:\,L(\tilde{\rho},\chi)=0}\hat{h}(\chi,\rho),
\end{equation}
where \(\vartheta(h)\) is the scaling operator on \(L^{2}_{0}(X_{\mathbb{Q}})\),
\(\rho\) ranges over zeros of Hecke \(L\)-functions \(L(s,\chi)\), and the
prime on the local integral denotes principal-value regularisation.
The \(\mathrm{GL}_{1}\) case with trivial character recovers
\(\xi(s)\): the non-trivial zeros of \(\zeta\) enter the absorption
trace distribution.
\end{theorem}

\begin{conjecture}[intertwiner comparing H-P\(^{\mathrm{Ar}}\) with Connes]
\label{v4-arith:hpvbconv:thm:connes-identification}
\ClaimStatusConjectured. Let \(H_{\mathrm{HP}}\) and
\(\Theta_{\xi}=1/2+iH_{\mathrm{HP}}\) be as in
\Cref{v4-arith:rh:hyp:HPAr}, and let \(D_{\mathrm{BC}}\) denote the
Friedrichs Fock/BC Hamiltonian of
\Cref{v4-arith:hpvbconv:thm:L0-Friedrichs}. Let
\(\vartheta\) be the scaling operator of Connes on
\(L^{2}_{0}(X_{\mathbb{Q}})\). The required comparison is the existence
of an intertwiner
\begin{equation}
\label{v4-arith:hpvbconv:eq:intertwiner}
\pi:\;\mathcal{V}^{\mathrm{prim}}_{\mathrm{Heis},\mathrm{prim}}\;\longrightarrow\;L^{2}_{0}(X_{\mathbb{Q}},d\mu)
\end{equation}
defined on \(\mathcal{D}_{0}\) and extended by closure, such that
\begin{equation}
\label{v4-arith:hpvbconv:eq:intertwining}
\pi\circ \Theta_{\xi}\;=\;\vartheta_{0}\circ\pi
\end{equation}
where \(\vartheta_{0}\) is the infinitesimal generator of the scaling
action on the kernel of the idele-class scaling trace. The intertwiner
\(\pi\) is \(C_{\mathbb{Q}}\)-equivariant on the sph-finite sector and
identifies the positive Weil--Connes trace of \(\Theta_{\xi}\) with
the distributional trace in Connes' trace formula. Its restriction to
the zero-mode lattice must also explain how \(D_{\mathrm{BC}}\), whose
partition trace is \(\zeta(\beta)\), maps to the absorption trace
without postulating the zero divisor.
\end{conjecture}

\begin{remark}[open intertwiner]
The sph-finite sector of \(\mathcal{V}^{\mathrm{prim}}\) embeds by
restricted tensor product into a subspace of \(L^{2}([GL_{2}(\mathbb{Q})]\backslash
GL_{2}(\mathbb{A}))\) compatible with the adele class space image
(Gelbart \emph{Automorphic Forms on Adele Groups} \S3.1). The
zero-mode lattice \(L\subset\mathcal{V}^{\mathrm{prim}}\) maps to the
finite-adele part of \(X_{\mathbb{Q}}^{0}\); the archimedean
Heisenberg factor \(\mathcal{H}^{(\infty)}\) maps to the
archimedean-fibre Fock of \(L^{2}_{0}(X_{\mathbb{Q}}^{0,\infty})\). The
scaling action of \(C_{\mathbb{Q}}\) on \(L^{2}_{0}(X_{\mathbb{Q}}^{0})\)
restricts on the sph-finite sector to the multiplicative action of
\(\mathbb{N}^{\times}\) on the Bost--Connes zero-mode lattice
(\Cref{v4-arith:thm:BC-acts-on-L}), plus the real-parameter
archimedean scaling on the Fock. The infinitesimal generator of
this action is the Bost--Connes/Fock Hamiltonian on the
chiral-algebra side, while the Connes trace formula sees the
absorption trace. Constructing an isometric intertwiner from the
former to the latter is exactly the determinant-trace problem.
\end{remark}

\begin{corollary}[H-P\(^{\mathrm{Ar}}\) and Connes have the same determinant-trace input]
\label{v4-arith:hpvbconv:cor:HPAr-connes-equivalence}
\ClaimStatusProvedHereConditional \emph{(conditional on the intertwiner of
\Cref{v4-arith:hpvbconv:thm:connes-identification})}. Under that
intertwiner, the trace-compatibility condition
\eqref{v4-arith:hpvbconv:eq:BC-trace-condition}, read as the
distributional equality of \Cref{v4-arith:rh:hyp:HPAr}, is equivalent
to Connes' absorption-trace hypothesis: the
\(C_{\mathbb{Q}}\)-scaling action on \(L^{2}_{0}(X_{\mathbb{Q}})\) has
zero distribution of \(\zeta\) with correct multiplicity.
\end{corollary}

\begin{proof}
Given the conjectural intertwiner, applying \(\pi\) transports the
regularised trace distribution of \(\Theta_{\xi}\) to the Connes
distributional trace. Conversely, on the sph-finite image of \(\pi\),
the Connes absorption-trace statement pulls back to the
Hilbert--P\'olya determinant and trace identities. The proof uses the
intertwiner as an input; it does not construct it.
\end{proof}

\begin{remark}[form of Connes' trace formula]
\label{v4-arith:hpvbconv:rem:connes-form}
Connes' trace formula \eqref{v4-arith:hpvbconv:eq:connes-trace-formula}
is proved in arXiv:math/9811068 Thm.~1. The identification of the
absorption trace with the zero distribution of \(\zeta\) with correct
multiplicity is the content of Connes' RH-equivalent statement; it is
\emph{not} a theorem. \Cref{v4-arith:hpvbconv:cor:HPAr-connes-equivalence}
places H-P\(^{\mathrm{Ar}}\) at the same level only after the
intertwiner of \Cref{v4-arith:hpvbconv:thm:connes-identification} is
assumed. That intertwiner is part of the RH-equivalent
determinant-trace problem. The chiral-algebra form replaces the
adele-class-space \(L^{2}\) setting by a Fock-space setting with
Sugawara and free-boson operators.
\end{remark}

\subsection{Connection to Deninger Motivic Cohomology}
\label{v4-arith:hpvbconv:deninger-connection}

\begin{definition}[Deninger axioms for motivic cohomology]
\label{v4-arith:hpvbconv:def:deninger-axioms}
\ClaimStatusDefinitional (Deninger Proc.~Symp.~Pure Math.~55.1, \S1).
A \emph{Deninger motivic cohomology} is a conjectural functor
\begin{equation}
H^{\bullet}_{M}(-,\mathbb{R}(\cdot)):\;\mathrm{Schemes}/\mathbb{Z}\;\longrightarrow\;\mathrm{Graded}\,\mathbb{R}\text{-vs}
\end{equation}
equipped with an integral Frobenius endomorphism \(\Phi\) satisfying:
\begin{itemize}
\item[(D1)] \emph{Lefschetz trace formula.} For every smooth
arithmetic scheme \(X\),
\(\zeta_{X}(s)=\prod_{i}\det_{\infty}(s\cdot\mathrm{id}-\Phi|H^{i}_{M}(X,\mathbb{R}(1)))^{(-1)^{i+1}}\).
\item[(D2)] \emph{Purity.} The eigenvalues of \(\Phi\) on
\(H^{i}_{M}(X,\mathbb{R}(1))\) all satisfy \(\mathrm{Re}(\rho)=i/2\) (or a
shifted variant).
\item[(D3)] \emph{Regularised determinant.} The determinant
\(\det_{\infty}\) is the \(\zeta\)-regularised determinant defined by
\(\log\det_{\infty}(A)=-\partial_{s}\mathrm{tr}(A^{-s})|_{s=0}\).
\end{itemize}
\end{definition}

\begin{proposition}[H-P\(^{\mathrm{Ar}}\) \(\Leftrightarrow\) Deninger formalism for \(\overline{\mathrm{Spec}\,\mathbb{Z}}\)]
\label{v4-arith:hpvbconv:prop:deninger-equivalence}
\ClaimStatusProvedHereConditional \emph{(conditional on existence of Deninger cohomology
for \(\overline{\mathrm{Spec}\,\mathbb{Z}}\) and on an explicit
Deninger--Fock intertwiner)}. Assuming
\Cref{v4-arith:hpvbconv:def:deninger-axioms} applies to
\(\overline{\mathrm{Spec}\,\mathbb{Z}}\) with \(X=\overline{\mathrm{Spec}\,\mathbb{Z}}\)
and \(H^{1}_{M}(\overline{\mathrm{Spec}\,\mathbb{Z}},\mathbb{R}(1))\)
non-zero, the Hilbert--P\'olya zero-divisor operator \(\Theta_{\xi}\) of
\Cref{v4-arith:rh:hyp:HPAr} is identified with
\(\Phi|H^{1}_{M}(\overline{\mathrm{Spec}\,\mathbb{Z}},\mathbb{R}(1))^{+}\)
(the \(+\) superscript denotes the primary
\(\mathrm{Gal}(\overline{\mathbb{Q}}/\mathbb{Q})\)-invariant part), and
H-P\(^{\mathrm{Ar}}\) holds iff (D1), (D2), (D3) hold.
\end{proposition}

\begin{proof}
Under the Deninger axioms, the meromorphic Tate zeta surface is a
cohomological superdeterminant: the even cohomology accounts for the
Tate poles at \(0\) and \(1\), while the primary \(H^{1}\) determinant
is the entire factor \(\xi_{\mathbb R}(s)\). The ordinary determinant
target on the primary summand is
\[
\xi_{\mathbb R}(s)=
\det_{\infty}(s\cdot\mathrm{id}-\Phi|H^{1}_{M}(-,\mathbb{R}(1))^{+}).
\]
The regularised determinant identity fixes the spectral divisor of the
primary summand:
\[
\mathrm{div}\,\det_{\infty}(s\cdot\mathrm{id}-\Phi|H^{1}_{M}(-,\mathbb{R}(1))^{+})
=\mathrm{div}\,\xi_{\mathbb R}(s).
\]
By (D3), the trace of \(\Phi\) is a distributional zero trace in the
admissible Weil--Connes sense, not an ordinary heat trace. The
identification is canonical only after composing the intertwiner of
\Cref{v4-arith:hpvbconv:thm:connes-identification} with the
conjectural Deninger--Connes dictionary of Leichtnam.
\end{proof}

\begin{remark}[the three formalisms in alignment]
\label{v4-arith:hpvbconv:rem:three-formalisms}
The three comparison formulations of H-P\(^{\mathrm{Ar}}\):
\begin{enumerate}
\item \emph{Fock/BC domain realisation}: the Friedrichs Hamiltonian
of \Cref{v4-arith:hpvbconv:thm:L0-Friedrichs} supplies a self-adjoint
Euler-side operator, but the passage to \(\Theta_{\xi}\) is the
determinant-trace realisation.
\item \emph{Connes adele class space}: the scaling generator
\(\vartheta_{0}\) on \(L^{2}_{0}(X_{\mathbb{Q}})\) has the zero
distribution as its absorption trace; this is Connes'
RH-equivalent reformulation of the trace formula of
arXiv:math/9811068 Thm.~1, not a formal consequence of it.
\item \emph{Deninger motivic cohomology}: the Frobenius \(\Phi\) on
\(H^{1}_{M}(\overline{\mathrm{Spec}\,\mathbb{Z}},\mathbb{R}(1))^{+}\)
has determinant divisor \(\mathrm{div}\,\xi_{\mathbb R}\)
(Deninger 1994 axioms).
\end{enumerate}
The equivalences among these formulations require the Connes/Fock and
Deninger/Fock intertwiners. Without those intertwiners there is a
comparison diagram, not a proof that the three operator models are the
same object.
\end{remark}

\subsection{Necessary and Sufficient Conditions}
\label{v4-arith:hpvbconv:necessary-sufficient}

\begin{theorem}[H-P\(^{\mathrm{Ar}}\) is the determinant-trace realisation]
\label{v4-arith:hpvbconv:thm:HPAr-equivalence}
\ClaimStatusProvedHere. The following are equivalent after the
explicit Connes/Fock and Deninger/Fock intertwiners named above are
included as hypotheses:
\begin{enumerate}
\item[(H-P\(^{\mathrm{Ar}}\))] The arithmetic Hilbert--P\'olya
hypothesis of \Cref{v4-arith:rh:hyp:HPAr}.
\item[(SP)] \emph{Determinant-trace condition:} there exists the
self-adjoint operator \(A\), the normal operator
\(\Theta_{\xi}=1/2+iA\), the determinant divisor, and the positive
Weil--Connes trace required in \Cref{v4-arith:rh:hyp:HPAr}.
\item[(Connes)] \emph{Connes trace realisation:} the
\(C_{\mathbb{Q}}\)-scaling action on \(L^{2}_{0}(X_{\mathbb{Q}})\) has
regularised trace distribution equal to the non-trivial zero
distribution of \(\zeta\), with multiplicity.
\item[(Deninger)] \emph{Deninger regularised-determinant condition:}
the formal identity
\(\xi_{\mathbb R}(s)=
\det_{\infty}(s\cdot\mathrm{id}-\Phi|H^{1}_{M,\mathbb{R}(1),+})\)
holds with \(\Phi\) a self-adjoint operator on a Hilbert space matching
\(\mathcal{V}^{\mathrm{prim}}_{\mathrm{Heis},\mathrm{prim}}\) via an
explicit intertwiner.
\end{enumerate}
\end{theorem}

\begin{proof}
(H-P\(^{\mathrm{Ar}}\)) \(\Leftrightarrow\) (SP) is definitional after
\Cref{v4-arith:rh:hyp:HPAr}: SP is the statement that the
self-adjoint operator, determinant divisor, and positive trace
distribution exist. The Friedrichs theorem supplies a
separate Fock/BC domain model; it is evidence for the analytic
framework but is not identified with \(H_{\mathrm{HP}}\) without an
intertwiner.

(SP) \(\Leftrightarrow\) (Connes) is
\Cref{v4-arith:hpvbconv:cor:HPAr-connes-equivalence}, conditional on
the Connes/Fock intertwiner. (SP) \(\Leftrightarrow\) (Deninger) is
\Cref{v4-arith:hpvbconv:prop:deninger-equivalence}, conditional on
the Deninger/Fock intertwiner and on the Deninger cohomology. Under
these named hypotheses, transitivity closes the comparison.
\end{proof}

\begin{corollary}[H-P\(^{\mathrm{Ar}}\) reduces to determinant and trace]
\label{v4-arith:hpvbconv:cor:HPAr-single-residual}
\ClaimStatusProvedHere. H-P\(^{\mathrm{Ar}}\) reduces to the
determinant-trace condition (SP). The Friedrichs/KMS side fixes the
natural Bost--Connes weight \(\lambda(m)=\log m\) and gives the Euler
partition trace; it does not give the zero trace. The required
assertion is the construction of \(H_{\mathrm{HP}}\),
\(\Theta_{\xi}\), and the positive Weil--Connes trace whose
distribution is the non-trivial-zero distribution of
\(\xi_{\mathbb R}\). This is the Hilbert--P\'olya content itself, not
a routine weight choice
(\Cref{v4-arith:hpvbconv:rem:hilbert-polya-content}).
\end{corollary}

\begin{proof}
Direct from \Cref{v4-arith:hpvbconv:thm:HPAr-equivalence}. The
trace-compatibility condition encodes the location of non-trivial
zeros of \(\xi_{\mathbb R}\); the Bost--Connes weight gives the Euler
side, while the zero side is the determinant divisor together with the
positive Weil--Connes trace distribution.
\end{proof}

\begin{remark}[Hilbert--P\'olya factor]
\label{v4-arith:hpvbconv:rem:hilbert-polya-content}
The Friedrichs construction of
\Cref{v4-arith:hpvbconv:thm:L0-Friedrichs} proves only the domain
theorem. Given \(\{\lambda(m)\}\geq 0\), the Friedrichs extension is
unique, self-adjoint, and non-negative. H-P\(^{\mathrm{Ar}}\) is the
additional determinant-trace assertion: one must construct a zero
operator with
\(\mathrm{div}\,\det_{\infty}(s-\Theta_{\xi})
=\mathrm{div}\,\xi_{\mathbb R}\) and trace distribution such that, in
symbolic wave-kernel notation,
\begin{equation*}
\mathrm{Tr}_{\mathrm{reg}}\bigl(e^{-t\Theta_{\xi}}\bigr)\;=\;\sum_{\rho}e^{-t\rho}\qquad(t>0),
\end{equation*}
where the equality is read after admissible Weil--Connes
regularisation and the sum is over non-trivial zeros of
\(\xi_{\mathbb R}\) counted with multiplicity. This is Hilbert's
original eighth problem, not a corollary of the Friedrichs step. The
determinacy comparison retains the domain theorem and leaves the
Hilbert--P\'olya determinant-trace realisation.
\end{remark}

\begin{remark}[Berry--Keating candidate]
\label{v4-arith:hpvbconv:rem:berry-keating}
Berry--Keating SIAM Review 41 (1999) proposed
\(H=\tfrac{1}{2}(xp+px)\) on \(L^{2}(\mathbb{R}_{>0},dx/x)\) as a candidate
Hilbert--P\'olya operator, restricted to a Planck-cell cutoff. This
operator has continuous spectrum on \(\mathbb{R}\) and the Bohr-Sommerfeld
quantisation of \(xp=E\) gives the correct leading asymptotic zero
density \(N(T)\sim (T/2\pi)\log(T/2\pi)\). In the present framework,
this candidate corresponds to a specific choice of
\(\{\lambda(m)\}_{m\in\mathbb{N}^{\times}}\): the prime-indexed weights
\(\lambda(p^{k})=k\log p\) (Bost--Connes monoid generators), combined
with the Fock archimedean spectrum \(\{i\,n_{i}\}\), formally give the
Berry--Keating spectrum structure. The Berry--Keating candidate is
\emph{not} rigorously known to have determinant divisor
\(\mathrm{div}\,\xi_{\mathbb R}\); it reproduces the density but not
the individual zero divisor. The rigorous
identification requires a regularisation matching
\eqref{v4-arith:hpvbconv:eq:BC-trace-condition}, which is the
H-P\(^{\mathrm{Ar}}\) content.
\end{remark}

\section{VB\(^{*}\) Converse: Structural Analysis}
\label{v4-arith:hpvbconv:VBstar-converse}

\subsection{The Moment Problem in Hilbert Space}
\label{v4-arith:hpvbconv:moment-problem}

\begin{theorem}[determinate vs.~indeterminate Hamburger moment problem]
\label{v4-arith:hpvbconv:thm:determinacy}
\ClaimStatusProvedElsewhere (Akhiezer 1965 \S2.4; Shohat--Tamarkin
1943 Ch.~III.5). Let \(\{\mu_{n}\}_{n\geq 0}\) be a sequence of real
numbers with Hankel matrices
\(H_{n}=(\mu_{i+j})_{0\leq i,j\leq n}\succeq 0\) for all \(n\). The
Hamburger moment problem is:
\begin{itemize}
\item \emph{Solvable}: there exists a positive measure \(\nu\) on
\(\mathbb{R}\) with \(\mu_{n}=\int x^{n}\,d\nu(x)\).
\item \emph{Determinate}: the solution \(\nu\) is unique.
\item \emph{Indeterminate}: multiple distinct positive measures share
the moment sequence.
\end{itemize}
Carleman's sufficient criterion for determinacy:
\(\sum_{n\geq 1}\mu_{2n}^{-1/(2n)}=+\infty\).
Non-sufficient examples (Stieltjes 1894): the lognormal moment
sequence \(\mu_{n}=e^{n^{2}/2}\) is indeterminate.
\end{theorem}

\begin{theorem}[Simon 1998: moment problem as Jacobi matrix]
\label{v4-arith:hpvbconv:thm:simon-jacobi}
\ClaimStatusProvedElsewhere (Simon Adv.~Math.~137 (1998) \S5). A
Hamburger moment problem with moments \(\{\mu_{n}\}_{n\geq 0}\) is
equivalent to the spectral theory of a Jacobi matrix
\[
J=\begin{pmatrix}b_{0}&a_{0}&0&\cdots\\ a_{0}&b_{1}&a_{1}&\cdots\\
0&a_{1}&b_{2}&\cdots\\ \vdots&\vdots&\vdots&\ddots\end{pmatrix}
\]
on \(\ell^{2}(\mathbb{Z}_{\geq 0})\), with \(a_{n}>0\) and \(b_{n}\in\mathbb{R}\)
determined recursively by the moment sequence via the
Gram--Schmidt-orthonormalised polynomial basis
\(\{p_{n}(x)\}_{n\geq 0}\) with \(p_{n}(x)=x^{n}+\text{lower}\) and
\(\int p_{m}p_{n}\,d\nu=\delta_{mn}\). The moment problem is
determinate iff \(J\) is essentially self-adjoint on \(c_{00}\), and
indeterminate iff \(J\) has deficiency indices \((1,1)\).
\end{theorem}

\begin{proposition}[VB\(^{*}\) operator-positivity vs.~moment positivity: the gap]
\label{v4-arith:hpvbconv:prop:VBstar-vs-moment-gap}
\ClaimStatusProvedHere. Let \(\nu\) be the symmetrised zero measure of
\Cref{v4-arith:vb:def:zero-measure}, with moments
\(\mu_{n}=\int x^{n}\,d\nu\). The following implications hold after the
stated Fock-space spectral identification:
\begin{enumerate}
\item VB\(^{*}\) (operator positivity: \(T_{n}\succeq 0\,\forall n\)) \(\Rightarrow\)
Li positivity (\(\lambda_{n}\geq 0\,\forall n\)).
\item Li positivity \(\Leftrightarrow\) RH \(\Leftrightarrow\) \(\nu\)
positive measure on the critical-line parameter set.
\item If \(\nu\) is determinate and the associated \(L_{0}^{\mathrm{Ar}}\)
has spectral support in \([1/2,\infty)\), then VB\(^{*}\) follows.
\item Positivity of the moments alone, without determinacy and this
support identification, does not imply VB\(^{*}\) for a general
unbounded moment problem.
\end{enumerate}
\end{proposition}

\begin{proof}
(1): trace of positive-semidefinite is non-negative; hence
\(\mathrm{tr}(T_{n})=n\lambda_{n}\geq 0\) under VB\(^{*}\).

(2): Li's criterion (Li 1997) and the moment reformulation of
\Cref{v4-arith:vb:thm:moment-equivalence}(2)\(\Leftrightarrow\)(4).

(3): under determinacy and the stated support condition,
\(L_{0}^{\mathrm{Ar}}\) has a unique spectral measure on the Fock
closure, and \(T_{n}\) diagonalises with
eigenvalues \(m^{n}-(m-1)^{n}\geq 0\) for \(m\geq 1/2\), hence
\(T_{n}\succeq 0\).

(4): in an indeterminate Hamburger problem there are distinct positive
measures with the same moments. They give distinct self-adjoint
realisations of the same formal Jacobi data. Moment positivity therefore
does not select the Fock operator, nor does it force the support
condition needed for \(m^{n}-(m-1)^{n}\geq0\). Akhiezer--Glazman
Vol.~II Ch.~VI.3 provides the classical model.
\end{proof}

\begin{corollary}[the converse input is spectral regularity]
\label{v4-arith:hpvbconv:cor:missing-is-determinacy}
\ClaimStatusProvedHere. The converse implication from zero-placement to
VB\(^{*}\) is the spectral regularity of the symmetrised zero moment
problem: positivity of the regularised moment functional, Hamburger
determinacy, and the Fock-space support identification placing
\(L_{0}^{\mathrm{Ar}}\) in \([1/2,\infty)\). Without these three
conditions, zero-placement does not force VB\(^{*}\).
\end{corollary}

\begin{proof}
Direct from
\Cref{v4-arith:hpvbconv:prop:VBstar-vs-moment-gap}(3)--(4): determinacy
is the hinge.
\end{proof}

\begin{proposition}[Carleman construction remains open]
\label{v4-arith:hpvbconv:prop:carleman-insufficient}
\ClaimStatusConjectured. Carleman's criterion
\(\sum_{n\geq 1}\mu_{2n}^{-1/(2n)}=+\infty\) is not yet verified for
the symmetrised zero moment sequence. The heuristic growth estimates
below are inconclusive; they do not justify claiming that Carleman's
criterion fails.
\end{proposition}

\begin{proof}
Under RH, \(\rho=1/2+i\gamma_{\rho}\) with \(\gamma_{\rho}\in\mathbb{R}\).
The symmetrised measure has support
\(\mu_{n}=\sum_{\rho}(\tfrac{1}{4}+\gamma_{\rho}^{2})^{n}\) regularised.
Using Riemann--von Mangoldt \(N(T)\sim (T/2\pi)\log(T/2\pi)\), the
\(N\)-th zero has \(\gamma_{N}\sim 2\pi N/\log N\). Bost--Connes
regularisation gives
\begin{equation*}
\mu_{2n}\;\asymp\;\sum_{N}\Bigl(\frac{2\pi N}{\log N}\Bigr)^{4n}\cdot(\text{regularisation
factor})\;\asymp\;(2\pi N_{\max})^{4n}/(\log N_{\max})^{4n}
\end{equation*}
where \(N_{\max}\) is the cutoff. If the regularised limit had the
growth \(\mu_{2n}\asymp C^{n}n!^{2}\) for some \(C>0\), then
\[
\mu_{2n}^{-1/(2n)}\sim (C^{n}n!^{2})^{-1/(2n)}
\asymp \frac{1}{\sqrt{C}\,n}.
\]
The harmonic comparison diverges, so this heuristic would support
Carleman's sufficient criterion rather than refute it. The required
input is a positive moment measure with sharp two-sided
regularised-moment asymptotics. Until those are proved, the Carleman
construction remains open.
\end{proof}

\begin{remark}[Stieltjes' example and the primary-Fock case]
\label{v4-arith:hpvbconv:rem:stieltjes}
Stieltjes 1894 gave an explicit indeterminate moment sequence:
\(\mu_{n}=e^{n^{2}/2}\). Riemann--von Mangoldt alone does not give a
two-sided growth law for the regularised zero moments.
The heuristic of \Cref{v4-arith:hpvbconv:prop:carleman-insufficient}
posits growth \(\mu_{2n}\asymp C^{n}n!^{2}\), which is between polynomial
and lognormal, as the plausible asymptotic until a rigorous construction is supplied.
Whether this intermediate rate is determinate is the
open question. It is classically known (Akhiezer 1965 \S2.4) that
\(\mu_{n}\asymp n!^{2}\) is on the boundary of Carleman's criterion;
determinacy depends on finer asymptotic data.
\end{remark}

\subsection{Hermite--Biehler + de Branges}
\label{v4-arith:hpvbconv:hermite-biehler}

\begin{definition}[Hermite--Biehler function]
\label{v4-arith:hpvbconv:def:hermite-biehler}
\ClaimStatusDefinitional (de Branges 1968 Ch.~II). An entire function
\(E(z)\) of finite exponential type is \emph{Hermite--Biehler} if
\(|E(z)|>|E^{*}(z)|\) for all \(z\) in the open upper half plane
\(\mathrm{Im}\,z>0\), where \(E^{*}(z)=\overline{E(\overline{z})}\).
The \emph{de Branges space} \(\mathcal{H}(E)\) is the reproducing-kernel
Hilbert space of entire functions \(f\) with \(f/E,f^{*}/E\in H^{2}(\mathbb{C}_{+})\)
(upper-half-plane Hardy space) and reproducing kernel
\begin{equation}
K_{E}(w,z)\;=\;\frac{\overline{E(w)}E(z)-\overline{E^{*}(w)}E^{*}(z)}{2\pi i(\overline{w}-z)}.
\end{equation}
\end{definition}

\begin{theorem}[de Branges structure theorem]
\label{v4-arith:hpvbconv:thm:de-branges-structure}
\ClaimStatusProvedElsewhere (de Branges 1968 Thm.~22). Every
reproducing-kernel Hilbert space of entire functions \(\mathcal{H}\)
satisfying:
\begin{itemize}
\item \(f,g\in\mathcal{H}\Rightarrow\overline{g(\overline{z})}\in\mathcal{H}\)
(complex conjugate closure);
\item for every non-real \(w\), the point evaluation \(f\mapsto f(w)\) is
continuous;
\item for every non-real \(w\),
\(f\mapsto(f-\tilde{f})/(z-w)\) is bounded (where \(\tilde{f}\) denotes a
specific orthogonal projection);
\end{itemize}
is of the form \(\mathcal{H}(E)\) for some Hermite--Biehler \(E\).
\end{theorem}

\begin{theorem}[Lagarias: RH via reproducing kernel]
\label{v4-arith:hpvbconv:thm:de-branges-1994}
\ClaimStatusProvedElsewhere (Lagarias Acta Arith.~89 (1999),
217--234). Let
\[
E_{\xi}(z)=\xi_{\mathbb R}(\tfrac{1}{2}-iz)+
\xi_{\mathbb R}'(\tfrac{1}{2}-iz).
\]
Then RH \(\Leftrightarrow\) \(E_{\xi}\) is Hermite--Biehler
\(\Leftrightarrow\) the reproducing kernel
\(K_{E_{\xi}}(w,z)\) is positive-semidefinite.
\end{theorem}

\begin{proposition}[VB\(^{*}\) converse via Hermite--Biehler refinement]
\label{v4-arith:hpvbconv:prop:VBstar-de-branges}
\ClaimStatusProvedHereConditional. Under RH and the Fock-space
identification of the de Branges model, the following are equivalent:
\begin{enumerate}
\item VB\(^{*}\) holds (operator-positivity \(T_{n}\succeq 0\) for all \(n\)).
\item The de Branges space \(\mathcal{H}(E_{\xi})\) is
\emph{operator-positive}: for every continuous linear functional
\(L:\mathcal{H}(E_{\xi})\to\mathbb{C}\) of the form
\(L(f)=\int f(x)\,d\sigma(x)\) with \(\sigma\) a signed measure, the
induced bilinear form
\begin{equation}
B_{L}(f,g)\;=\;\int f(x)\overline{g(x)}\,d\sigma(x)
\end{equation}
is positive-semidefinite on the kernel of \(L_{0}^{\mathrm{Ar}}-1/2\).
\item The spectral-regularity condition of
\Cref{v4-arith:hpvbconv:cor:missing-is-determinacy} holds.
\end{enumerate}
\end{proposition}

\begin{proof}[Sketch]
(1) \(\Leftrightarrow\) (3): \Cref{v4-arith:hpvbconv:cor:missing-is-determinacy}.

(3) \(\Leftrightarrow\) (2): de Branges 1968 Ch.~III.24 relates
determinacy of a moment problem to uniqueness of the positive
reproducing-kernel extension containing the polynomials densely. The
``operator-positive'' condition (2) is the corresponding positivity and
support condition after transport to the primary-Fock image.

(1) \(\Leftrightarrow\) (2): follows by transitivity (1) \(\Leftrightarrow\) (3) \(\Leftrightarrow\) (2)
without invoking the conjectural intertwiner. Alternatively, when the Connes/Fock
intertwiner of
\Cref{v4-arith:hpvbconv:thm:connes-identification} is assumed, the primary-Fock
\(\mathcal{V}^{\mathrm{prim}}_{\mathrm{Heis},\mathrm{prim}}\) embeds
into \(\mathcal{H}(E_{\xi})\) via that intertwiner composed with the
de Branges Mellin transform of Connes 1999; under this embedding,
\(T_{n}\) corresponds to the multiplication operator by
\(x^{n}-(x-1)^{n}\) on \(\mathcal{H}(E_{\xi})\), making (2) equivalent to
operator-positivity of that multiplication operator.
\end{proof}

\subsection{The B\'aez-Duarte Refinement}
\label{v4-arith:hpvbconv:baez-duarte}

\begin{theorem}[Nyman--Beurling 1950/1955]
\label{v4-arith:hpvbconv:thm:nyman-beurling}
\ClaimStatusProvedElsewhere (Nyman 1950; Beurling 1955). Let
\(\rho(x)=\{1/x\}-\lfloor 1/x\rfloor\) be the fractional part function.
Let \(M_{0}\subset L^{2}((0,1))\) be the closed linear span of
\(\{\rho(\theta x):\theta\in(0,1]\}\). Then RH \(\Leftrightarrow\)
\(\mathbf{1}_{(0,1)}\in M_{0}\).
\end{theorem}

\begin{theorem}[B\'aez-Duarte 2003]
\label{v4-arith:hpvbconv:thm:baez-duarte}
\ClaimStatusProvedElsewhere (B\'aez-Duarte Adv.~Math.~190 (2005),
315--320; arXiv:math/0202141). Let \(\rho_{n}(x)=\{n/x\}-n\{1/x\}\) for
\(n\geq 2\). Then RH \(\Leftrightarrow\) the functions
\(\{\rho_{n}\}_{n\geq 2}\) span a dense subspace of \(L^{2}((0,1),dx)\).
Let
\begin{equation}
\epsilon_{N}^{\mathrm{BD}}\;:=\;\mathrm{dist}_{L^{2}}\bigl(\mathbf{1},\mathrm{span}\{\rho_{n}\}_{2\leq n\leq
N}\bigr)
\end{equation}
be the finite-dimensional distance. Any sharp asymptotic for
\(\epsilon_{N}^{\mathrm{BD}}\) is an additional quantitative statement.
\end{theorem}

\begin{proposition}[conditional B\'aez-Duarte comparison]
\label{v4-arith:hpvbconv:prop:VBstar-BD-rate}
\ClaimStatusProvedHereConditional. Assume that the Mellin transform
identifies the finite B\'aez-Duarte distance with the corresponding
Petersson distance in the primary arithmetic Heisenberg Fock, with the
normalising constant fixed. Then the B\'aez-Duarte condition used here
is exactly the Petersson-matched form of the VB\(^{*}\) spectral
regularity condition.
\end{proposition}

\begin{proof}[Sketch]
The Nyman--Beurling \(M_{0}\)-space embeds into the primary arithmetic
Heisenberg Fock via the Mellin transform
\begin{equation*}
\rho_{n}(x)\;\longmapsto\;n^{-s}\Bigl(\zeta(s)-\frac{n}{s-1}\Bigr)
\end{equation*}
(Burnol 2000, Thm.~1.2). The distance
\(\epsilon_{N}^{\mathrm{BD}}\) corresponds to the
Petersson-Mellin distance from \(\mathbf{1}_{(0,1)}\) (the primary
vacuum) to the zero-mode-lattice span of
\(\{\rho_{n}\}_{2\leq n\leq N}\) (the Bost--Connes orbit of the vacuum).
Under the stated identification, positivity of \(T_{n}\) is the
Hankel-positivity and determinacy condition for the same projected
vacuum vector. Parseval's identity then transfers the Petersson
normalisation to the \(L^{2}\)-distance
\(\epsilon_{N}^{\mathrm{BD}}\). No unconditional estimate for
\(\epsilon_{N}^{\mathrm{BD}}\) is asserted here.
\end{proof}

\begin{theorem}[VB\(^{*}\) converse and spectral regularity]
\label{v4-arith:hpvbconv:thm:VBstar-converse-spectral-regularity}
\label{v4-arith:hpvbconv:thm:VBstar-de-branges-BD}
\ClaimStatusProvedHereConditional. Assume P2, P3, H-P\(^{\mathrm{Ar}}\),
and the Fock-space identifications stated in
\Cref{v4-arith:hpvbconv:prop:VBstar-de-branges} and
\Cref{v4-arith:hpvbconv:prop:VBstar-BD-rate}.
The following are equivalent:
\begin{enumerate}
\item[(VB\(^{*}\)c)] The VB\(^{*}\) converse: zero-placement \(\Rightarrow\) VB\(^{*}\).
\item[(Reg)] positivity of the regularised zero-moment functional,
Hamburger-determinacy of the symmetrised zero moment sequence of
\Cref{v4-arith:vb:def:zero-measure}, and the Fock support condition of
\Cref{v4-arith:hpvbconv:cor:missing-is-determinacy}.
\item[(HB+)] \(E_{\xi}\) is Hermite--Biehler \emph{and} the associated
de Branges space \(\mathcal{H}(E_{\xi})\) is operator-positive in the
sense of \Cref{v4-arith:hpvbconv:prop:VBstar-de-branges}(2).
\item[(BD+)] RH holds \emph{and} the B\'aez-Duarte decay rate is
\(\epsilon_{N}^{\mathrm{BD}}\asymp 1/\sqrt{\log N}\) with the Petersson-matched
constant.
\end{enumerate}
This theorem is a structural comparison. It does not assert that RH
alone implies any of (Reg), (HB+), or (BD+).
\end{theorem}

\begin{proof}
(VB\(^{*}\)c) \(\Leftrightarrow\) (Reg):
\Cref{v4-arith:hpvbconv:cor:missing-is-determinacy}: RH plus
spectral regularity is precisely the content of the VB\(^{*}\) converse
in this framework.

(Reg) \(\Leftrightarrow\) (HB+):
\Cref{v4-arith:hpvbconv:prop:VBstar-de-branges}.

(Reg) \(\Leftrightarrow\) (BD+):
\Cref{v4-arith:hpvbconv:prop:VBstar-BD-rate}.

Carleman's sufficient criterion has not been verified for the
symmetrised zero moment sequence
(\Cref{v4-arith:hpvbconv:prop:carleman-insufficient}). Abstract
indeterminate moment problems show that unbounded trace data do not
formally determine operator positivity. They do not prove an
independence theorem for the zeta zero sequence.
\end{proof}

\begin{remark}[the converse does not help the forward direction]
\label{v4-arith:hpvbconv:rem:converse-does-not-help}
\Cref{v4-arith:hpvbconv:thm:VBstar-converse-spectral-regularity}
and \Cref{v4-arith:hpvbconv:cor:three-strictly-finer} are
comparisons at the VB\(^{*}\) converse level: each of (Reg),
(HB+), (BD+) is a form of the same spectral-regularity input for
zero-placement \(\Rightarrow\) VB\(^{*}\). The forward direction
VB\(^{*}\Rightarrow\)zero-placement is the separate content of
\Cref{v4-arith:rh:thm:VB-implies-F-RH}, which is a conditional
theorem under P2, P3, H-P\(^{\mathrm{Ar}}\), and a positive
Bost--Connes regularised trace. The conditional comparison for the
converse does \emph{not} close the forward implication, and does
not discharge any of the conditional hypotheses of the forward
direction. In particular, the sufficient RH path of
\Cref{v4-arith:rh:thm:sufficient-path} requires VB\(^{*}\) main
direction plus H-P\(^{\mathrm{Ar}}\); the comparison for
the converse identifies the analytic content of VB\(^{*}\) but does
not reduce the forward requirement.
\end{remark}

\subsection{The Spectral-Regularity Content}
\label{v4-arith:hpvbconv:strictly-finer}

\begin{corollary}[three equivalent spectral-regularity conditions]
\label{v4-arith:hpvbconv:cor:three-strictly-finer}
\ClaimStatusProvedHereConditional. Under P2, P3, H-P\(^{\mathrm{Ar}}\)
and the stated Fock-space identifications, the following three
conditions are equivalent forms of the VB\(^{*}\) converse:
\begin{enumerate}
\item VB\(^{*}\) (arithmetic Virasoro operator positivity for all \(n\)).
\item The spectral-regularity condition (Reg).
\item B\'aez-Duarte sharp decay \(\epsilon_{N}^{\mathrm{BD}}\asymp
1/\sqrt{\log N}\) with the Petersson constant (BD+).
\end{enumerate}
Condition (HB+) (Hermite--Biehler + operator-positivity of the de
Branges space) is a fourth equivalent form.
\end{corollary}

\begin{proof}
Direct from \Cref{v4-arith:hpvbconv:thm:VBstar-converse-spectral-regularity}.
\end{proof}

\begin{remark}[the content of the spectral-regularity gap]
\label{v4-arith:hpvbconv:rem:strict-content}
The four conditions (VB\(^{*}\), Reg, HB+, BD+) express the same
additional spectral regularity after the stated identifications. Classical
abstract moment-problem theory (Akhiezer 1965 \S 2.4, Simon 1998
\S 5) constructs unbounded self-adjoint operators for which moment data
do not determine the operator. Hence the VB\(^{*}\) converse requires an
input not contained in RH as a bare zero-placement statement. The
content of this input is, in each form:
\begin{itemize}
\item VB\(^{*}\): operator-positivity of \(T_{n}=(L_{0}^{\mathrm{Ar}})^{n}-(L_{0}^{\mathrm{Ar}}-\mathrm{id})^{n}\)
on the primary arithmetic Heisenberg Fock.
\item Reg: positivity, uniqueness, and support of the spectral measure
determined by the symmetrised zero moment sequence.
\item HB+: operator-positivity of the de Branges reproducing kernel
\(\mathcal{H}(E_{\xi})\) extending beyond Hermite--Biehler.
\item BD+: a precise \(L^{2}\) convergence rate of Beurling-sequence
approximations of \(\mathbf{1}_{(0,1)}\).
\end{itemize}
These four conditions are equivalent among themselves only under the
hypotheses of
\Cref{v4-arith:hpvbconv:thm:VBstar-converse-spectral-regularity}. A
proof of RH which did not also identify the Fock spectral measure would
not close VB\(^{*}\) in this argument.
\end{remark}

\section{Relation Between H-P\(^{\mathrm{Ar}}\) and VB\(^{*}\) Converse}
\label{v4-arith:hpvbconv:relation}

\begin{theorem}[H-P\(^{\mathrm{Ar}}\) and VB\(^{*}\) reduce to RH plus regularity]
\label{v4-arith:hpvbconv:thm:HPAr-implies-reduction}
\ClaimStatusProvedHereConditional. Under H-P\(^{\mathrm{Ar}}\) and the
Fock-space support identification, VB\(^{*}\) is equivalent to RH
\emph{plus} the spectral-regularity condition of
\Cref{v4-arith:hpvbconv:cor:missing-is-determinacy}. Without
H-P\(^{\mathrm{Ar}}\), the implication RH \(\Rightarrow\) VB\(^{*}\) is
not meaningful, since VB\(^{*}\) presupposes a self-adjoint
\(L_{0}^{\mathrm{Ar}}\).
\end{theorem}

\begin{proof}
Direct from
\Cref{v4-arith:hpvbconv:thm:VBstar-converse-spectral-regularity}
combined with
\Cref{v4-arith:hpvbconv:thm:HPAr-equivalence}: once H-P\(^{\mathrm{Ar}}\)
is in place, \(L_{0}^{\mathrm{Ar}}\) has a well-defined spectral
measure, and VB\(^{*}\) \(\Leftrightarrow\) RH plus spectral regularity
reduces to the moment-problem equivalence with the support condition.
\end{proof}

\begin{proposition}[H-P\(^{\mathrm{Ar}}\) and regularity are separate data]
\label{v4-arith:hpvbconv:prop:HPAr-det-independent}
\ClaimStatusProvedHere. H-P\(^{\mathrm{Ar}}\) (existence of a
self-adjoint \(L_{0}^{\mathrm{Ar}}\) with prescribed determinant
divisor and trace distribution) and the
spectral-regularity condition (positivity, determinacy, and support of
the spectral measure) are separate pieces of data:
\begin{itemize}
\item There exist indeterminate moment sequences with positive Hankel
structure (Stieltjes' \(e^{n^{2}/2}\) example); these admit multiple
self-adjoint realisations. H-P\(^{\mathrm{Ar}}\) asserts existence of
one realisation with the right trace; regularity asserts uniqueness and
support properties of the realisation.
\item Conversely, a determinate moment sequence does not automatically
admit a self-adjoint realisation that matches a specific target
trace; the target trace (the Bost--Connes trace of
\Cref{v4-arith:hpvbconv:eq:BC-trace-condition}) is an additional
constraint on the spectrum.
\end{itemize}
\end{proposition}

\begin{proof}
Akhiezer 1965 \S 2.4 and Simon 1998 Thm.~5.3 show that a moment
sequence may admit several positive measures. Existence of one
self-adjoint realisation and uniqueness of the realisation are therefore
distinct assertions. The prescribed Bost--Connes trace is a further
constraint.
\end{proof}

\begin{corollary}[separation of the two analytic inputs]
\label{v4-arith:hpvbconv:cor:residuals-independent}
\ClaimStatusProvedHere. The two analytic inputs to VB\(^{*}\) in the
sufficient path of \Cref{v4-arith:rh:thm:sufficient-path}, H-P\(^{\mathrm{Ar}}\)
and the VB\(^{*}\) converse, are not the same assertion. Hence both are
required for the VB\(^{*}\) input, and neither is redundant.
\end{corollary}

\begin{proof}
\Cref{v4-arith:hpvbconv:prop:HPAr-det-independent}.
\end{proof}

\begin{remark}[the refined sufficient path]
\label{v4-arith:hpvbconv:rem:refined-sufficient-path}
The sufficient path to RH of \Cref{v4-arith:rh:thm:sufficient-path}
has four inputs: arithmetic Positselski, Koszul--Petersson
compatibility, H-P\(^{\mathrm{Ar}}\), and VB\(^{*}\). The last two
split as:
\begin{itemize}
\item H-P\(^{\mathrm{Ar}}\): existence of a self-adjoint operator
obtained from \(L_{0}^{\mathrm{Ar}}\) by the required intertwiner, whose
normal zero-divisor operator has determinant divisor
\(\mathrm{div}\,\xi_{\mathbb R}\) and whose
Weil--Connes trace is the corresponding zero distribution. Equivalent
forms: Connes adele-class trace realisation; Deninger motivic
regularised determinant with purity.
\item VB\(^{*}\): RH \emph{and} the spectral-regularity condition for the
symmetrised zero moment sequence. Equivalent forms: Hermite--Biehler
operator-positivity; B\'aez-Duarte sharp rate after the Petersson
normalisation has been fixed.
\end{itemize}
Each of the two inputs is a single spectral condition.  The
Connes--Deninger, de Branges, and Nyman--Beurling forms are equivalent
languages for those conditions.
\end{remark}

\section{Analytic Conditions}
\label{v4-arith:hpvbconv:residual}

The two analytic conditions have the following forms:

\begin{enumerate}
\item \textbf{(H-P\(^{\mathrm{Ar}}\), Friedrichs construction).}
Construct the determinant-trace operator and the trace-compatible embedding
for which condition~\eqref{v4-arith:hpvbconv:eq:BC-trace-condition}
holds on the Friedrichs extension of
\Cref{v4-arith:hpvbconv:thm:L0-Friedrichs}. Equivalent forms:
\begin{itemize}
\item Connes: the \(C_{\mathbb{Q}}\)-scaling action on
\(L^{2}_{0}(X_{\mathbb{Q}})\) has absorption trace equal to the zero
distribution of \(\zeta\), with correct multiplicity.
\item Deninger: \(\Phi|H^{1}_{M}(\overline{\mathrm{Spec}\,\mathbb{Z}},\mathbb{R}(1))^{+}\)
has determinant divisor \(\mathrm{div}\,\xi_{\mathbb R}\) with correct
multiplicity.
\end{itemize}
\item \textbf{(VB\(^{*}\) converse, regularity construction).} Establish
positivity, Hamburger determinacy, and the support identification for
the symmetrised zero moment sequence \(\{\mu_{n}\}_{n\geq 0}\) of
\Cref{v4-arith:vb:def:zero-measure}.
Equivalent forms:
\begin{itemize}
\item de Branges: operator-positivity of the reproducing kernel on
\(\mathcal{H}(E_{\xi})\) beyond Hermite--Biehler.
\item B\'aez-Duarte: sharp Petersson-normalised rate for
\(\epsilon_{N}^{\mathrm{BD}}\).
\end{itemize}
\item \textbf{(Combined VB\(^{*}\) input).} Both (1) and (2)
together supply the VB\(^{*}\) input in the sufficient path of
Vol~IV. Neither alone suffices: (1) without (2) gives a specific
self-adjoint operator but leaves \(T_{n}\)-positivity open; (2) without
(1) gives abstract moment regularity without a Fock-space
operator.
\end{enumerate}

\begin{remark}[local form]
\label{v4-arith:hpvbconv:rem:gain-summary}
The Hilbert--P\'olya condition is existence together with determinant
divisor and trace distribution.  The VB\(^{*}\) converse is positivity,
Hamburger determinacy, and support for the symmetrised zero moments.
The Connes--Deninger and de Branges--B\'aez-Duarte forms identify these
conditions in classical analytic languages; they do not make either
condition a formal consequence of the other.
\end{remark}

\section{Spectral Boundary}
\label{v4-arith:hpvbconv:summary}

\begin{enumerate}
\item \textbf{H-P\(^{\mathrm{Ar}}\) is determinant-trace realisation}
(\Cref{v4-arith:hpvbconv:thm:HPAr-equivalence},
\Cref{v4-arith:hpvbconv:cor:HPAr-single-residual}): the Friedrichs
extension of \Cref{v4-arith:hpvbconv:thm:L0-Friedrichs} closes
the domain problem, but the determinant divisor
\(\mathrm{div}\,\xi_{\mathbb R}\) and the matching Weil--Connes zero
trace are the genuine Hilbert--P\'olya content
(\Cref{v4-arith:hpvbconv:rem:hilbert-polya-content}). The
formulation is equivalent to Connes' adele-class trace realisation and
to Deninger's motivic cohomology
regularised-determinant condition.
\item \textbf{The VB\(^{*}\) converse is equivalent to spectral
regularity of the symmetrised zero-moment sequence}
(\Cref{v4-arith:hpvbconv:thm:VBstar-converse-spectral-regularity}),
under the stated Fock-space identifications. This regularity is not
supplied by zero placement alone.
The converse is an equivalence at the converse level only and
does \emph{not} close the forward direction VB\(^{*}\Rightarrow\)zero-placement
(\Cref{v4-arith:hpvbconv:rem:converse-does-not-help}).
Carleman's criterion remains an open possible construction
(\Cref{v4-arith:hpvbconv:prop:carleman-insufficient}); abstract
indeterminate moment problems show why this extra hypothesis
cannot be suppressed. Equivalent forms: de Branges
reproducing-kernel operator-positivity beyond Hermite--Biehler;
B\'aez-Duarte sharp decay rate with Petersson-matched constant.
\item \textbf{The two analytic inputs are separate}
(\Cref{v4-arith:hpvbconv:prop:HPAr-det-independent}): H-P\(^{\mathrm{Ar}}\)
asserts existence of the operator; regularity asserts positivity,
uniqueness, and support of the spectral measure. Both are required for
the VB\(^{*}\) input in the sufficient path; the other inputs of that path
are not supplied here.
\item \textbf{Berry--Keating \(H=\tfrac{1}{2}(xp+px)\) candidate} is a candidate choice
of Bost--Connes zero-mode weights \(\lambda(p^{k})=k\log p\), reproducing
the leading asymptotic zero density but not the individual zeros;
rigorous identification requires the regularisation matching
\eqref{v4-arith:hpvbconv:eq:BC-trace-condition}.
\item \textbf{Stieltjes' example of an indeterminate moment sequence}
(\(\mu_{n}=e^{n^{2}/2}\)) shows that the unbounded case admits genuine
indeterminacy; heuristic estimates suggest the symmetrised zero moment
sequence has intermediate growth \(\asymp C^{n}n!^{2}\) on the boundary of
Carleman, but this estimate is not yet rigorous
(\Cref{v4-arith:hpvbconv:rem:stieltjes}).
\item \textbf{The refined sufficient path to RH} has four inputs
(\Cref{v4-arith:hpvbconv:rem:refined-sufficient-path}): arithmetic
Positselski, Koszul--Petersson compatibility, H-P\(^{\mathrm{Ar}}\)
with explicit Friedrichs domain and trace-compatibility, and VB\(^{*}\)
as RH plus spectral regularity. The last two are each reducible to
single spectral conditions with classical comparison formulations.
\item \textbf{Independence of the two analytic inputs}: H-P\(^{\mathrm{Ar}}\)
does not imply the VB\(^{*}\) converse, and the VB\(^{*}\) converse
does not construct the Hilbert--P\'olya operator.  The four classical
languages, Connes, Deninger, de Branges, and B\'aez-Duarte, identify
the same two conditions from independent sides.
\end{enumerate}
